% =============================================================================
% AAAI-27 AISI submission draft (anonymous) --- assembled from section drafts.
% Uses the OFFICIAL AAAI-27 author kit (aaai2027.sty / aaai2027.bst), copied
% from ./authorkit/AuthorKit27/. Compile: pdflatex + bibtex + pdflatex x2.
% AISI keyword (entered in OpenReview, not here):
%   "AISI: Computational Social Science and Humanities"
%
% Working title below; alternatives proposed by section writers:
%   - "The Camera Did Not Move: Pose Forensics for Tracing Archival
%      Photographs into Wartime Propaganda Print"
%   - "Exposure Identity: Separating Reprints from Re-Shots in Wartime
%      Propaganda Archives with 3D Foundation-Model Camera Pose"
% =============================================================================
\documentclass[letterpaper]{article} % DO NOT CHANGE THIS
\usepackage[submission]{aaai2027}  % DO NOT CHANGE THIS
\usepackage[hyphens]{url}  % DO NOT CHANGE THIS
\usepackage{graphicx} % DO NOT CHANGE THIS
\urlstyle{rm} % DO NOT CHANGE THIS
\def\UrlFont{\rm}  % DO NOT CHANGE THIS
\usepackage{natbib}  % DO NOT CHANGE THIS AND DO NOT ADD ANY OPTIONS TO IT
\usepackage{caption} % DO NOT CHANGE THIS AND DO NOT ADD ANY OPTIONS TO IT
\frenchspacing  % DO NOT CHANGE THIS
% Recommended by the kit; permitted packages only (no hyperref/geometry/etc.).
\usepackage{booktabs}
\usepackage{amsmath}
\usepackage{amssymb}
\usepackage{xcolor}
\pdfinfo{
/TemplateVersion (2027.1)
}
\setcounter{secnumdepth}{0}

% --- Draft-only TODO macro: strip all \todo{}, remove \usepackage{xcolor},
% and delete all draft comment blocks (incl. the trailing reproducibility
% block and refs.bib comments) before submission. --------------------------
\newcommand{\todo}[1]{\textcolor{red}{\textbf{[TODO: #1]}}}

\title{Same Photograph or Same Scene? Camera-Pose Forensics for Tracing Archival Photographs into Wartime Propaganda Print}
\author{Anonymous submission}
\affiliations{Anonymous submission}

\begin{document}

\maketitle

\begin{abstract}
Wartime propaganda magazines reprinted archival photographs after cropping,
retouching, and recaptioning them; the record of that editorial labor survives
only as pairwise identities between archival exposures and printed
reproductions. Recovering those identities at scale is a verification problem,
not a retrieval one: the historian must decide whether a print derives from
the \emph{same exposure} as an archival photograph, or from a second
photograph of the same scene taken moments apart---the error that turns
provenance tools into misattribution machines. We present a four-stage
pipeline: self-supervised cross-domain retrieval, dense-correspondence
filtering, homography inlier-ratio gating, and a new forensic signal---the
relative camera pose inferred by a feed-forward 3D foundation model (VGGT).
Re-editions of one exposure behave as pictures of a picture and yield
near-zero inferred pose displacement; same-scene negatives do not. With all
thresholds frozen on a development shard, the pose gate raises held-out
verification precision from 83.7\% (at 100\% recall) to 90.2\% at 98.4\%
recall, conditional on the candidate set. We release NCR-Match, 2{,}636
expert-verified photograph--print pairs with exposure-identity versus
same-scene labels drawn from the North China Railway Archive, a CC~BY~4.0
occupation-era collection, and evaluate cross-system transfer to an opposing
propaganda apparatus and historian-in-the-loop utility.
\end{abstract}

\section{Introduction}
\label{sec:intro}

In March 1942 the \emph{Jinchaji Pictorial}, the flagship photographic organ
of the Chinese Communist base areas, printed a photograph by Sha Fei of the
peasant Liu Hanxing enlisting in the Eighth Route Army, under the page title
``Model Families: Parents Send Their Children, Wives Send Their Boys.'' The
same exposure was reprinted at least three more times over the next four
years; by 1945 Liu's name had been removed and the caption read ``common
people,'' and by 1946 he had become simply ``a youth''
\cite{du2024probing,du2025metaphor}. No single page records this
trajectory. The editorial labor of wartime propaganda---cropping, retouching,
splitting, recaptioning---is invisible in the printed object and undocumented
by the institutions that performed it: as Du observes, such archives preserve
the photographs but not the decisions \cite{du2025metaphor}. Propaganda
systems operate precisely by
recirculating a stable stock of authorized images under shifting textual
frames \cite{yurchak2005everything}, so the record of that work survives only
as thousands of pairwise identities between an archival exposure and a printed
reproduction---each individually minor, collectively the only evidence of how
the apparatus curated reality, what Du calls ``the evidence of the small''
\cite{du2026evidence}. The scale defeats close reading alone
\cite{arnold2019distant}: the North China Railway Archive, the photographic
record of the Japanese-run railway company at the center of occupation
publicity in North China \cite{kushner2006thought,earhart2008certain}, holds
roughly 39{,}000 photographs (1939--1945), an unknown fraction of which
resurfaced in the coeval propaganda pictorial \emph{North China}
(\emph{Hokushi}, 1939--1943). Figure~\ref{fig:teaser} shows the pair type at
issue.

\begin{figure}[t]
\centering
\fbox{\parbox{0.94\columnwidth}{\centering\vspace{2em}
\todo{Figure 1: one verified photo--halftone reprint pair vs.\ one
same-scene-seconds-apart negative, with inferred camera frusta overlaid ---
render from the pose-validity visualization batch (W11, report section 6.5)}
\vspace{2em}}}
\caption{A verified reprint pair (one exposure, re-edited into print) versus
a same-scene negative exposed moments apart, with the inferred cameras
overlaid.}
\label{fig:teaser}
\end{figure}

Computer vision appears to be the obvious remedy, and published work has shown
that learned embeddings can surface photograph--print candidates across the
severe domain gap of halftone screening, contrast loss, and cropping
\cite{du2024probing}. But retrieval answers the wrong question. Retrieval
finds look-alikes; the historian's claim---\emph{this print is that
negative}---is a claim of \emph{exposure identity}. Press photographers worked
in bursts, and archives are dense with frames taken seconds apart from nearly
the same position; a system that cannot separate ``same photograph, edited''
from ``same scene, next frame'' does not merely lose precision points, it
manufactures confident misattributions of photographer, place, and date. The
failure mode is not hypothetical: the second round of our benchmark's
verification protocol excluded a nontrivial fraction of first-round
candidates as exactly such same-scene negatives (see Dataset Construction).
The distinction itself is classical---image phylogeny separates near-duplicates
(one exposure, a chain of transformations) from semantically similar images of
the same scene \cite{dias2013forests,chum2007scalable}---yet modern
copy-detection benchmarks are born-digital and treat same-scene shots as
negatives by construction \cite{douze2021disc,pizzi2022sscd,wang2023ndec}, and
no public benchmark poses this boundary on archival photo-to-print material.
Where \citet{du2024probing} demonstrated proof-of-concept circulation
forensics with retrieval alone, this paper contributes the verification
decision retrieval cannot supply, delivered at benchmark scale and released
with the submission. Nor is
the problem antiquarian: recaptioned reuse of authentic photographs is the
dominant mechanism of contemporary out-of-context misinformation, for whose
modern form this community has built detectors and benchmarks
\cite{aneja2023cosmos,qi2024sniffer}. A wartime propaganda corpus is the same
mechanism with ground truth---a closed, fully attributable natural experiment
in institutional image reuse, eighty years before the feed.

We answer the exposure-identity question with a four-stage funnel
(Figure~\ref{fig:pipeline}). A self-supervised fusion head over frozen vision
backbones (DINOv3 \cite{simeoni2025dinov3}, EfficientNetV2, Swin) retrieves
top-ranked print candidates for each archival photograph; ASpanFormer
\cite{chen2022aspanformer} extracts dense correspondences and discards pairs
with too few; a seeded robust homography fit (MAGSAC++) gates on the inlier
ratio and aligns the photograph onto the print; and a final judge inspects a
signal new to
provenance analysis: the relative camera pose inferred by the feed-forward 3D
foundation model VGGT \cite{wang2025vggt}. The intuition is geometric. A
genuine reprint is a picture of a picture: after alignment the pair carries no
parallax, and the model infers essentially the same camera for both frames. A
second exposure of the same scene, however similar, retains parallax that
surfaces as rotation and lateral translation of the inferred camera. We
convert this into a component-weighted pose-displacement score---strict on
rotation and $x$/$y$ translation, lenient on the depth and field-of-view terms
that cropping and rescanning perturb---and threshold it. While feed-forward 3D
foundation models have been used to verify whether two images depict the same
surface or place
\cite{cai2023doppelgangers,xiangli2025doppelgangerspp,maggio2026vggtslam},
and copy detection has relied on 2D appearance descriptors with same-scene
shots treated as negatives by construction
\cite{douze2021disc,pizzi2022sscd,wang2023ndec}, we are---to our
knowledge---the first to use the relative camera pose inferred by such models
as a provenance signal separating re-editions of the same exposure (crops,
edits, halftone reprints) from distinct exposures of the same scene. On the
held-out validation shard of our benchmark, with every threshold frozen on the
development shard, the full system attains verification precision 90.2\% and
recall 98.4\% ($F_1$ 94.1), conditional on the candidate set; the pose gate
contributes $+6.5$ points of precision over inlier-ratio verification alone
(83.7\% $\rightarrow$ 90.2\%) at a cost of 1.6 points of recall.

\noindent This paper makes five contributions:
\begin{itemize}
\item \textbf{Task.} We formalize \emph{exposure-identity vs.\ same-scene
  discrimination} for archival photo-to-print tracing, grounding it in the
  phylogeny dichotomy between near-duplicate and semantically similar images
  \cite{dias2013forests} and the IND/NIND (identical vs.\ non-identical
  near-duplicate) taxonomy
  \cite{chum2007scalable,jindaapiraksa2013california}, and argue that it is
  the decision on which archival provenance claims---and their modern
  out-of-context analogues---rest.
\item \textbf{Method.} Inferred relative camera pose as a provenance signal,
  with the component-weighted decision rule above and ablations isolating it
  against verification without pose, classical pose recovered for free from
  the already-computed homography and from essential-matrix decomposition,
  MASt3R pose \cite{leroy2024mast3r}, and a learned classifier over identical
  inputs \todo{Table~\ref{tab:verification} rows B2, B4--B10: P/R/F1 + PR-AUC
  per variant with bootstrap CIs and McNemar tests vs.\ the full system ---
  CPU ablations over cached manifests plus MASt3R run W6}.
\item \textbf{Benchmark.} NCR-Match: 2{,}636 verified
  photograph--print pairs linking the North China Railway Archive
  ($\sim$39{,}000 photographs, 1939--1945) to the magazine \emph{North China}
  (1939--1943), embedded in a working corpus of 42{,}773 cropped images,
  curated with a two-round verification protocol in which same-scene
  near-misses are excluded as labeled hard negatives. Archive images and
  metadata are CC~BY~4.0, so the benchmark is deposited with this submission
  with attribution to the archive. To our knowledge it is the first public
  benchmark with exposure-identity versus same-scene labels on cross-domain
  (photograph vs.\ halftone print) archival material. Labels are
  expert-adjudicated under a written codebook \todo{inter-annotator agreement:
  Cohen's $\kappa$ from the dual-annotation campaign, report section 6.6}.
\item \textbf{Cross-system transfer.} With all thresholds frozen, the pipeline
  transfers from Japanese occupation propaganda to the opposing Communist
  base-area system: the \emph{Jinchaji Pictorial} / Papers of Sha Fei corpus
  (2{,}146 photographs, 1{,}122 pages, 4{,}390 crops, 230 annotated pairs)
  studied by prior retrieval-only work
  \cite{du2024probing}. \todo{Retrieval Recall@10/15 and mAP over the 230
  verified pairs + verification P/R/F1 conditional on hand-labeled filter
  survivors, thresholds frozen --- Jinchaji run W10. Jinchaji side releases
  annotations and identifiers only, no images.}
\item \textbf{Historian-in-the-loop evaluation and findings.} An annotation
  reliability study \todo{$\kappa$, adjudicated gold labels, `Unsure' as a
  first-class outcome}, a within-subjects time-saved study \todo{4--6
  historians, time-to-first-verified-republication and verified matches per
  30-minute session, pipeline-verified candidates vs.\ plain top-10 list},
  and new historical findings surfaced by the system \todo{2--3 case studies
  from the confirmed-match list, adjudicated by the domain expert}.
\end{itemize}

The system is built for a deployment path rather than a demonstration: the
anonymized code in the supplementary material, the benchmark, and the
annotation codebook are deposited with this submission, and the pipeline is
designed to slot into the verification workflows of the archives and research
groups that steward these collections. Concretely, the candidate-review
interface used in the historian study of the Experiments section is part of
the code release; the verified-pair layer is exported in a tabular,
identifier-keyed format suitable for ingestion by an archive's public
catalog; and the frozen-threshold transfer protocol demonstrated on the
Jinchaji corpus doubles as the documented onboarding recipe for applying the
system to a new collection. Its design is ethics-aware by
principle: \emph{the pipeline does not interpret; it retrieves}
\cite{du2026evidence}. Every accepted pair ships with its geometric evidence,
every threshold is published, ``Unsure'' is reported rather than silently
dropped, and the interpretive claim---what a recaptioned reprint
\emph{meant}---remains the historian's to make. We return to the dual-use
risks of automated provenance assertion, and to the contextualization owed to
imagery produced by occupation and propaganda apparatuses, in the Ethics
section.

\section{Related Work}
\label{sec:related}

\subsection{Content-Based Retrieval for Archives}
Deep metric learning made instance-level retrieval practical
\cite{schroff2015facenet}, and self-supervised contrastive objectives such as
SimCLR proved effective for learning transformation-invariant descriptors
\cite{chen2020simclr}. Transformer encoders sharpened global retrieval
features further \cite{dosovitskiy2021vit,elnouby2021training}; today,
self-supervised and image--text foundation encoders---DINOv2 and DINOv3
\cite{oquab2024dinov2,simeoni2025dinov3}, CLIP \cite{radford2021clip}, and
SigLIP~2 \cite{tschannen2025siglip2}---supply strong retrieval features off
the shelf. Matching across a change of medium or style is the harder,
cross-domain variant of the problem
\cite{shrivastava2011crossdomain,hu2023crossdomain}, and the
photograph-to-halftone gap studied here is an extreme instance of it. We
therefore treat retrieval as a component rather than a contribution: our fused
encoder exists to propose candidates, and we evaluate it against the
off-the-shelf encoders above.

\subsection{Copy Detection and the Same-Scene Boundary}
Near-duplicate and copy detection is the retrieval special case closest to our
task. Early benchmarks were small and gentle: Copydays covers 157 source
images under simple edits \cite{douze2009copydays}. DISC21 scaled the problem
to a million-image reference set under adversarial transformations, but its
queries and references remain born-digital \cite{douze2021disc}, and SSCD
distilled the resulting lessons into the standard copy-detection descriptor
\cite{pizzi2022sscd}. The frontier these efforts name is precisely our problem
class: the NDEC benchmark introduced hard negative pairs that are inherently
similar rather than edited copies, with an asymmetrical similarity objective
to handle them \cite{wang2023ndec}, and AnyPattern extends copy detection to
unseen edit patterns \cite{wang2026anypattern}. The distinction itself is
older. Near-identical image detection was posed by \citet{chum2007scalable},
and California-ND annotates a personal photo collection with exactly the split
we need: identical near-duplicates (variants of one exposure) versus
non-identical near-duplicates (different exposures of similar content)
\cite{jindaapiraksa2013california}. Most directly, the image-phylogeny
literature formalized the boundary over a decade ago: Dias, Rocha, and
Goldenstein reconstruct transformation trees over near-duplicates
\cite{dias2012phylogeny} and, in extending trees to forests, define the
\emph{semantically similar} images that must be kept out of a phylogeny---the
same scene captured at a ``different point in space and time''
\cite{dias2013forests}. Provenance analysis operationalizes phylogeny at scale
with 2D local features and appearance descriptors \cite{moreira2018provenance},
standardized by the NIST Media Forensics Challenge protocols \cite{guan2019mfc}
and recently rebuilt on vision transformers
\cite{zhang2024provenancevit,lai2025provenance}. Across this entire line, the
copy-versus-scene boundary is enforced by appearance alone; no system consults
the camera geometry that defines it. To our knowledge, ours is the first
provenance system to enforce the phylogeny boundary geometrically rather than
by appearance, and the first tested on cross-domain, photograph-to-print
archival material.

\subsection{Geometric Verification and 3D Foundation Models}
Confirming that a retrieved candidate corresponds to a query, rather than
merely resembling it, is classically the job of local features and robust
estimation: SIFT correspondences filtered by RANSAC
\cite{lowe2004sift,fischler1981ransac}. Detector-free transformer matchers
such as LoFTR \cite{sun2021loftr} and ASpanFormer \cite{chen2022aspanformer}
find dense correspondences under weak texture and heavy degradation, where
halftone dot screens defeat classical detectors; LightGlue
\cite{lindenberger2023lightglue} and RoMa
\cite{edstedt2024roma,edstedt2025romav2} push accuracy and efficiency further.
In our pipeline this stage is a learned form of the historian's side-by-side
comparison. A newer family of feed-forward 3D foundation models regresses
scene geometry directly: DUSt3R and MASt3R predict pointmaps and matches from
image pairs \cite{wang2024dust3r,leroy2024mast3r}, Reloc3r regresses relative
camera pose at scale \cite{dong2025reloc3r}, and VGGT jointly predicts
cameras, depth, and point tracks for arbitrary image sets \cite{wang2025vggt}.
Deciding per image pair whether a homography (no parallax) or a fundamental
matrix (genuine two-view parallax) better explains the correspondences is a
classical model-selection problem \cite{torr1998gric}; our pose gate is its
forensic descendant, replacing information criteria over sparse
correspondences with halftone-robust dense matching and the metric pose head
of a 3D foundation model. The nearest neighbors of that use are Doppelgangers
and Doppelgangers++, which learn to reject visually convincing matches between
\emph{different} surfaces to disambiguate structure from motion
\cite{cai2023doppelgangers,xiangli2025doppelgangerspp}; our decision has the
opposite polarity---both images genuinely show the same surface, and the
question is whether they are the same exposure of it. VGGT-SLAM~2.0 uses VGGT
internals as a training-free verifier of candidate loop closures
\cite{maggio2026vggtslam}: same-place verification inside SLAM, not
duplicate-versus-reshot forensics. Computational re-photography estimated the
pose delta between a historical photograph and a live viewfinder to guide a
photographer back to the original viewpoint \cite{bae2010rephotography}---
inter-photograph camera pose as a first-class quantity, but in an interactive
rather than forensic setting. Same-photo discrimination itself appears once
before, on 2D appearance: detecting whether two face crops derive from one
photograph, to expose bias in kinship benchmarks \cite{dawson2018samephoto}---
no geometry, and never extended to provenance. A related intuition, that a
photograph of a flat print or screen betrays itself geometrically, underlies
recapture detection in image-authenticity verification
\cite{vilesov2024deepfakes}; that family makes single-image attack decisions,
not pairwise provenance decisions. In sum, none of these lines uses inferred
relative camera pose to decide exposure identity; that use, argued in the
Introduction, is ours.

\subsection{Why Not Sensor Forensics?}
The natural forensic reflex for the same-exposure question is sensor-noise
attribution: PRNU fingerprints identify the source camera of a digital image
\cite{lukas2006prnu} and survive, weakly, into printed copies
\cite{goljan2008printed}. Three facts rule this channel out here. First,
wartime photographs were exposed on film; there is no digital sensor and hence
no PRNU to extract. Second, a device fingerprint identifies the camera, not
the exposure: two frames taken seconds apart with the same camera share it, so
it answers the wrong question. Third, halftone screening and generational
duplication destroy the high-frequency signal the method depends on, a
degradation documented in the printed-image study itself
\cite{goljan2008printed}. Geometry survives exactly where sensor noise dies:
the projective relation between two prints of one exposure is untouched by
screening.

\subsection{Historical Image Reuse and Cross-Domain Datasets}
Reuse of visual material predates photography: the cost of woodblocks and
copperplates led early modern printers to recycle blocks across editions, and
object detection plus visual search now recovers those block-sharing networks
\cite{lang2025finetuning,dutta2021chapbooks}; contrastive training likewise
clusters recurring printing ornaments across eighteenth-century books
\cite{wang2025ornamentreuse}. For photographic archives, Newspaper Navigator
established the page-level visual-extraction lineage from which our cropping
stage descends \cite{lee2020newspaper}, and transductive fine-tuning has been
applied to near-duplicate detection in scanned photo collections
\cite{net2023transductive}---the archival near-duplicate setting, addressed
with 2D descriptors and no geometry. Purpose-built historical corpora are
growing: EyCon offers roughly 130k digitized photographs and press pages from
the conflicts of 1890--1918 \cite{giardinetti2024eycon}; the Illustrated
London News dataset pairs some 72k wood engravings with multimodal search
\cite{smits2025iln}; and FORBIN releases a press-photography agency archive of
the early twentieth century with rich archival metadata and an explicit
image-reuse research agenda \cite{chelali2026forbin}. The closest
methodological neighbor is the EyCon group's halftone similarity work
\cite{aissi2025halftone}, which combines visual and textual similarity to
improve similarity search over halftone press photographs. Two differences
position our task against theirs: their goal is intra-corpus similarity
neighborhoods within a single digitized collection, whereas ours is open-set
matching across two collections---an archive of original photographs and a
printed magazine; and their output is a ranked neighborhood, whereas ours is a
verified pair-level decision about exposure identity. None of these corpora
carries verified pair ground truth linking a photograph to its printed
reproduction; NCR-Match fills that gap (see the contributions above). Finally,
\citet{du2024probing} demonstrated proof-of-concept circulation forensics on
WWII East Asian pictorials with an ensemble retrieval model, reporting 77.8\%
top-15 retrieval on the Jinchaji corpus, and Du traced the editorial cropping
and recaptioning of such photographs through close reading
\cite{du2025metaphor}. Building on that proof of concept, this paper
contributes what neither work possessed: a scaled, benchmarked, and released
matching system, and a verification signal that decides whether two images
share an exposure.

\subsection{Out-of-Context Misuse and Computational Propaganda}
The decision our system makes---same exposure or not---is the evidentiary
substrate of a problem this community already studies. Out-of-context misuse,
in which an authentic image recirculates under a false caption, is the
cheapest and most common form of visual misinformation; COSMOS detects it by
grounding captions against image regions \cite{aneja2023cosmos}, and SNIFFER
adds multimodal LLM-based explanation \cite{qi2024sniffer}. Media scholarship
documents the same mechanics across a century: iconic photographs are revived
and resignified as memes \cite{boudana2017icons}, and their online
circulations can be traced at the scale of hundreds of thousands of reposts
\cite{smits2023iconic}. A 1943 propaganda magazine recaptioning an archival
photograph is mechanically the same act as a modern out-of-context post---the
wartime corpus is a fully labeled natural experiment in institutional image
reuse, with the added difficulty that each ``repost'' passed through a
halftone press. Methodologically, we invert distant viewing
\cite{arnold2019distant}: rather than reading a large corpus for aggregate
patterns, Digital Historical Forensics uses scale to locate the singular
object and then reads it closely \cite{du2026evidence}, with the division of
labor between pipeline and historian fixed by the design principle stated in
the Introduction.

\section{Task and the NCR-Match Benchmark}
\label{sec:task}

\subsection{Exposure Identity vs.\ Same Scene}

Given a query archival photograph $q$ and a candidate reproduction $c$
cropped from a printed page, the task is to decide whether the pair is
\emph{exposure-identical}: whether both images derive from a single release
of the shutter---one negative---so that $c$ is related to $q$ by a chain of
reproduction transformations (cropping, rotation, mirroring, rescaling,
retouching, halftone screening, and generational re-photography of prints).
The contrasting case is a \emph{same-scene negative}: $q$ and $c$ depict the
same scene but originate from distinct exposures---adjacent frames on a roll,
or frames made moments apart from a nearby vantage point. The dichotomy
itself is established. Image phylogeny distinguishes near-duplicates, related
by a transformation chain from a common source, from ``semantically similar''
images captured by a camera at a different point in space and time
\cite{dias2012phylogeny,dias2013forests}, and near-duplicate retrieval
separates identical from non-identical near-duplicates (IND vs.\ NIND)
\cite{chum2007scalable,jindaapiraksa2013california}. NCR-Match
operationalizes this boundary as a labeling standard on cross-domain archival
material, where it carries evidentiary weight: for the provenance claim
``this archival photograph is the one printed on this page,'' a same-scene
negative is not a weaker match but a false one.

Archival practice makes such negatives dense rather than exotic.
Photographers on assignment exposed sequences of near-identical frames, and
picture editors selected among the variants, sometimes printing different
frames of one sequence in different issues \cite{du2025metaphor}; wartime
pictorials also recirculated the same negative across publications
\cite{du2024probing}. The confusable pairs in this corpus therefore arise
organically from production practice, unlike born-digital copy-detection
benchmarks, whose negatives are unrelated by construction and whose
same-scene confusables are rare or synthetic
\cite{douze2021disc,wang2023ndec}. On our development shard, 316 of 641
expert-labeled candidate pairs (49\%) are negatives, in large part same-scene
look-alikes (\todo{exact same-scene vs.\ unrelated split of labeled
negatives, from the negative-type audit of Shards 1--3}).
Table~\ref{tab:dataset} summarizes the benchmark and its shards.

\begin{table}[t]
\centering
\small
\setlength{\tabcolsep}{4pt}
\begin{tabular}{@{}lr@{}}
\toprule
\multicolumn{2}{@{}l}{\emph{Source material}}\\
Archive photographs (NCR Archive, 1939--1945) & $\sim$39{,}000 \\
\emph{North China} magazine issues (1939--1943) & \todo{issue count} \\
Working corpus (YOLOv11 crops, both collections) & 42{,}773 \\
\midrule
\multicolumn{2}{@{}l}{\emph{Exposure-identity verification (two rounds)}}\\
First-round candidate pairs & \todo{count} \\
Excluded in round two as same-scene & \todo{count} ($\sim$7\%) \\
Verified pairs (NCR-Match) & 2{,}636 \\
\midrule
\multicolumn{2}{@{}l}{\emph{Labeled candidate pairs (pos.\,/\,neg.\,/\,Unsure)}}\\
Shard 1 (development) & 325\,/\,316\,/\,\todo{n} \\
Shard 2 (validation, thresholds frozen) & 252\,/\,386\,/\,\todo{n} \\
Shard 3 (test, evaluated once) & \todo{n\,/\,n\,/\,n} \\
Inter-annotator agreement (Cohen's $\kappa$) & \todo{$\kappa$} \\
\midrule
\multicolumn{2}{@{}l}{\emph{Transfer domain (Jinchaji Pictorial\,/\,Papers of Sha Fei)}}\\
Photographs\,/\,magazine pages & 2{,}146\,/\,1{,}122 \\
Crops\,/\,annotated pairs & 4{,}390\,/\,230 \\
\bottomrule
\end{tabular}
\caption{NCR-Match at a glance. Verification labels exist for candidate pairs
that survive retrieval and the correspondence floor; verification metrics are
therefore conditional on the candidate set (see Experiments), and end-to-end
recall is reported separately as a survival analysis over the 2{,}636
verified pairs.}
\label{tab:dataset}
\end{table}

\subsection{Dataset Construction}

\textbf{Sources.} NCR-Match pairs photographs from the North China Railway
Archive---approximately 39{,}000 photographs taken in Japanese-occupied North
China between 1939 and 1945, digitized and published with metadata under
CC~BY~4.0 \cite{ncrarchive2019}\footnote{\url{https://codh.rois.ac.jp/north-china-railway/}}---with
reproductions in the Japanese-language pictorial \emph{North China}
(\emph{Hokushi}, 1939--1943). Magazine pages mix several photographs with
text; following the page-segmentation lineage of Newspaper Navigator
\cite{lee2020newspaper}, a YOLOv11 detector \cite{khanam2024yolov11} trained
on a small annotated page sample crops individual photographs. The working
corpus comprises 42{,}773 cropped images across the two collections;
crop-stage quality is \todo{detection precision/recall on a 100-page manually
audited sample --- YOLOv11 QC audit}.

\textbf{Two-round verification.} Candidate photo--print pairs were proposed
by embedding retrieval with geometric matching over the working corpus
(\todo{first-round candidate-pair count}) and verified by hand in two rounds:
round one screened candidates for visual identity; round two re-examined
every accepted pair specifically for the same-scene failure mode. About 7\%
of first-round candidates were identified in round two as same-scene
negatives and excluded, leaving 2{,}636 verified exposure-identity pairs,
each carrying magazine issue, date, and page metadata. The exclusion rate is
itself evidence for the task: even under careful single-pass review, roughly
one accepted candidate in fourteen was a different exposure of the same
scene. The excluded pairs are retained as labeled same-scene hard negatives
(\todo{confirm the round-two exclusion list ships in the release}).

\textbf{Labeled verification shards.} For threshold development and testing,
the working corpus is deterministically partitioned into six shards of
roughly 7{,}000 images each; three are labeled in this work. Within a shard,
each query photograph's ten highest-ranked print candidates are matched with
ASpanFormer \cite{chen2022aspanformer}, and pairs passing the correspondence
floor (roughly 650 per shard; see Method) receive human labels. Verification
labels thus exist only for pairs the funnel surfaces, and all verification
metrics are reported conditional on the candidate set (see Experiments).
Shard~1 is the development set---all thresholds and the pose-component
weighting were chosen here---with 325 positive and 316 negative pairs.
Shard~2 is held-out validation, evaluated once with everything frozen:
252\,/\,386. Shard~3 is an untouched test shard, labeled under the codebook
below and evaluated exactly once: \todo{Shard-3 positive/negative/Unsure
counts, from the Shard-3 labeling campaign}. Unsure labels (\todo{Unsure
counts per shard, from the dual-annotation campaign}) are reported, never
silently dropped.

\textbf{Transfer domain.} A second, disjoint corpus supports cross-domain
evaluation: the \emph{Jinchaji Pictorial} and the Papers of Sha Fei (2{,}146
photographs, 1{,}122 magazine pages, 4{,}390 crops), with 230 annotated
photo--print pairs from \citet{du2024probing}. The two corpora come from
opposing propaganda
systems---Japanese occupation publishing and Communist base-area
publishing---with different printing technology and degradation profiles.

\subsection{Annotation Protocol and Reliability}

Labels are governed by a written codebook whose rulings fix the two genuinely
contested cases: adjacent frames from one photographing sequence are
same-scene negatives, however similar; and disjoint crops from a single
negative are positives---exposure identity does not require overlapping image
content. Each candidate pair is labeled independently by two annotators, a
domain-expert historian of wartime East Asian print media and a trained
annotator, with \emph{Unsure} as a first-class label; disagreements are
adjudicated into gold labels. Inter-annotator agreement is \todo{Cohen's
$\kappa$ over the dual, independently labeled Shard 1--3 candidate pairs
($\approx$1{,}950), from the relabeling campaign}. Every verified positive
additionally carries a trace type---halftone-only, crop, rotation/mirror,
retouch, or recaption---with distribution \todo{trace-type distribution over
the 2{,}636 verified positives, expert taxonomy tagging}; per-type pipeline
pass rates appear in the Experiments section. The recaption class---
near-identical pixels under a changed caption---is detectable as an image
match, but its editorial significance is invisible to an image-only system
(see Limitations).

\subsection{Release, Licensing, and Known Biases}

NCR-Match is deposited with the supplementary material, together with
anonymized code, rather than promised: image crops, pair labels, shard
assignments, trace annotations, the codebook, and per-pair metadata. The
North China Railway Archive publishes its photographs and metadata under
CC~BY~4.0, and the release attributes all NCR imagery to the archive
accordingly. For the Jinchaji domain we release annotations and image
identifiers only, no image data. A datasheet \cite{gebru2021datasheets}
documents collection, the two-round verification protocol, annotation, and
known biases. Chief among these: the source archive is the product of
Japanese corporate and military photography---it records what an occupation
apparatus chose to photograph---so both subject matter and the trace-type
distribution reflect that selection. We state this as a documented limitation
on generalization claims rather than treating the corpus as a neutral sample
of wartime photography (see Ethics).

\section{Method}
\label{sec:method}

\begin{figure}[t]
\centering
\fbox{\parbox{0.94\columnwidth}{\centering\vspace{2em}
\todo{pipeline figure: four-stage funnel with one photo--print pair as
running example (archival photograph, retrieved print crop, ASpanFormer
correspondences, homography-aligned overlay, VGGT pose deltas), annotated
with per-stage survivor counts from the survival audit (W5/W11)}
\vspace{2em}}}
\caption{Overview of the four-stage verification funnel.}
\label{fig:pipeline}
\end{figure}

Figure~\ref{fig:pipeline} gives an overview. Given an archival photograph
(source) and a corpus of cropped print images (targets), the pipeline is a
four-stage funnel: (1)~\emph{fusion retrieval} proposes, for each source, its
top-$K$ most similar targets; (2)~\emph{detector-free matching} establishes
dense correspondences and discards pairs below a correspondence floor;
(3)~a \emph{geometric-consistency gate} requires the correspondences to be
explainable by a single planar projective map; (4)~a \emph{pose-based
decision layer} labels each surviving pair \emph{same exposure} or
\emph{same scene, different exposure}. Stages are ordered by cost, and each
stage writes every measured signal, decision, and machine-readable rejection
reason to an append-only manifest, so all threshold-level ablations replay
from cached artifacts without GPU recomputation; every randomized estimator
is seeded, and the released code reproduces the reported decisions
bit-for-bit. Because the funnel restricts what later stages see, every metric
downstream is reported against a named denominator (see Experiments).

\subsection{Stage 1: Fusion Retrieval}

The retrieval encoder follows the frozen-ensemble design validated on
archival photo--print retrieval by \citet{du2024probing}, with the ViT
backbone upgraded to DINOv3. Each image is resized to $224{\times}224$ and
rendered as three-channel grayscale before ImageNet normalization: across
the photograph/halftone gap, tonal structure, not color, is the informative
channel. Three frozen backbones encode the image---DINOv3 ViT-L/16
\citep[1{,}024-D;][]{simeoni2025dinov3}, EfficientNetV2-M
\citep[1{,}000-D;][]{tan2021efficientnetv2}, and Swin-T
\citep[1{,}000-D;][]{liu2021swin}---and their concatenated 3{,}024-D output
is projected by a two-layer MLP head ($3{,}024 \!\to\! 2{,}000$, ReLU,
$\to\! 1{,}000$). Only the head is trained, with SimCLR-style
self-supervision \citep{chen2020simclr} on the unlabeled 42{,}773-crop
working corpus: two augmented views (random crops and rotations) per image,
contrastive temperature $0.05$, batches of 16 images (32 views), 50 epochs,
SGD with learning rate $10^{-5}$, momentum $0.9$, weight decay $10^{-3}$,
fixed seed. No pair labels are used at any stage of training. Because the
corpus mixes both domains, the head learns halftone-invariant features;
because it is also the corpus later searched, the setup is transductive,
which we disclose and control \todo{transductive control: retrain the
fusion head excluding all Shard-3 images and show Shard-3 verification F1
unchanged within the bootstrap CI --- run W9}. Embeddings are
$\ell_2$-normalized and compared by cosine similarity; the $K{=}10$
highest-scoring targets per source become candidate pairs \todo{retrieval
Recall@K (K=5--30) and mAP over verified pairs, per encoder incl.\
off-the-shelf baselines, Table~\ref{tab:retrieval} --- run W4 from the cached
\texttt{feature\_cache.npz}}.

\subsection{Stage 2: Detector-Free Matching and Correspondence Floor}

Halftone screening, retouching, and paper degradation defeat detector-based
local features, so we match with ASpanFormer \citep{chen2022aspanformer}
(pretrained outdoor weights; grayscale; long side resized to 1{,}024\,px,
aspect preserved), which regresses dense correspondences without keypoint
detection. For each candidate pair we record the raw correspondence set and
its epipolar-consistent subset under a fundamental-matrix fit with seeded
RANSAC \citep[\texttt{FM\_RANSAC}, 1\,px reprojection threshold,
\texttt{cv2.setRNGSeed} before every robust
fit;][]{fischler1981ransac,bradski2000opencv}. A pair advances only if its
correspondence count clears a floor $\tau_{\mathrm{corr}}$
\todo{BLOCKING, resolve before any other table row is run: state
$\tau_{\mathrm{corr}}$ for the Shard-1/2 runs (working log records
50/75/100 in different places), resolve whether the floor applies to raw
correspondences or to epipolar inliers (the code gates on the filtered
count, the log narrative says raw), and verify that the resolved definition
regenerates exactly the 641/638 labeled candidate sets of Shards~1--2 that
produced rows B2/B3; floor swept 0--200 in
Table~\ref{tab:verification} row~B1, CPU-only from
\texttt{aspan\_all\_manifest.jsonl} --- run W1}. The floor is a compute
gate, not a match decision: it keeps the expensive pose stage off hopeless
pairs, and because every attempted pair---including failures---is logged,
the floor can be swept post hoc without re-running the matcher.

\subsection{Stage 3: Geometric-Consistency Gate}

From the stored raw correspondences we estimate a source$\to$target
homography $H$ with seeded USAC\_MAGSAC \citep[5\,px
threshold;][]{barath2020magsac} and define the inlier ratio $\rho$ as the
fraction of raw correspondences consistent with $H$. If two images are
re-editions of one exposure, the mapping between them is, to good
approximation, a single planar projective transform---cropping, rescanning,
rotation, and page geometry compose into one homography---so $\rho$ is
high. Correspondences that scatter off a single homography indicate either
scene-level similarity with parallax or unstable matching. The gate accepts
pairs with $\rho \ge 0.65$, derived on the development shard and frozen
thereafter. A deliberately lenient sanity pre-pass (finite, non-singular
$H$; minimal correspondence, inlier, and warped-overlap requirements) skips
the pose stage only for egregiously degenerate alignments, which are
recorded as rejections.

\subsection{Stage 4: Pose-Based Exposure-Identity Decision}

The homography then serves a second purpose: the source is warped onto the
target canvas, and the aligned pair is passed jointly to VGGT
\citep{wang2025vggt} \todo{pin the exact VGGT checkpoint, commit hash, and
inference resolution in the reproducibility checklist}, a feed-forward 3D
foundation model that regresses for each frame a nine-dimensional camera
encoding $[\mathbf{t}, \mathbf{q}, \mathbf{f}]$: translation
$\mathbf{t}\in\mathbb{R}^3$ in the model's normalized scene units, an
orientation quaternion $\mathbf{q}$, and vertical/horizontal field of view
$\mathbf{f}\in\mathbb{R}^2$. From the two frames' encodings we take four
deltas: the rotation angle
$\delta_{\mathrm{rot}} = 2\arccos\lvert\langle\hat{\mathbf{q}}_s,
\hat{\mathbf{q}}_t\rangle\rvert$ (unit-normalized quaternions; the absolute
dot product resolves the $\mathbf{q}\equiv-\mathbf{q}$ ambiguity), the
lateral shift $\delta_{xy} = \lVert\mathbf{t}^{(xy)}_s -
\mathbf{t}^{(xy)}_t\rVert_2$, the depth shift $\delta_{z} = \lvert t^{(z)}_s
- t^{(z)}_t\rvert$, and the zoom shift $\delta_{\mathrm{fov}} =
\lVert\mathbf{f}_s - \mathbf{f}_t\rVert_2$. Each is scaled, capped, and
weighted into a single score
\begin{equation}
S \;=\!\! \sum_{c\,\in\,\{\mathrm{rot},\,xy,\,z,\,\mathrm{fov}\}}\!\!
w_c \,\min\!\left(\frac{\delta_c}{s_c},\, \gamma\right),
\label{eq:pose-score}
\end{equation}
with weights $(w_{\mathrm{rot}}, w_{xy}, w_{z}, w_{\mathrm{fov}}) =
(1, 1, 0.25, 0.25)$, scales $(s_{\mathrm{rot}}, s_{xy}, s_{z},
s_{\mathrm{fov}}) = (2^{\circ}, 0.020, 0.100, 0.100)$, and per-term cap
$\gamma = 3$. Here $\delta_{\mathrm{rot}}$ is converted to degrees before
scaling, and the translation deltas are measured in VGGT's normalized scene
units, so each ratio $\delta_c / s_c$ is dimensionless. The asymmetry is the
design: re-orientation and lateral
viewpoint change are physical evidence that the camera moved---a different
exposure---whereas apparent zoom and depth shifts routinely arise from
cropping, enlargement, and reproduction even in true copies, so those terms
contribute weakly and saturate \todo{empirical justification: per-component
violin/ROC plots of $\delta_{\mathrm{rot}}, \delta_{xy}, \delta_{z},
\delta_{\mathrm{fov}}$ on exact-copy vs.\ same-scene labeled pairs (CPU,
from \texttt{vggt\_judged\_manifest.jsonl}), Figure~\ref{fig:pose-components};
component ablations rotation+xy-only vs.\ fov+z-only,
Table~\ref{tab:verification} rows B4/B5}. A pair is declared the same
exposure iff
\begin{equation}
\rho \ge 0.65 \;\;\wedge\;\; S \le 2.13,
\label{eq:decision}
\end{equation}
both thresholds derived once on the development shard, frozen for all other
data, and committed as named constants in the released code. VGGT's global
register-token cosine similarity is recorded for every pair but is
deliberately \emph{excluded} from the rule: its distributions on exact
copies and same-scene negatives overlap heavily and a gate on it adds no
separation---a negative finding we report as such
\todo{Table~\ref{tab:verification}: global-similarity-gate baseline row B6
(cosine $\ge 0.90$) plus the distribution-overlap figure, CPU from cached
manifests}.

\subsection{Why Pose Can Judge a Picture of a Picture}

An objection must be met head-on: an aligned exact copy is a photograph of
a photograph---planar, degenerate two-view geometry, far from the
multi-view scenes a 3D foundation model is trained on. Can its pose
estimates be trusted there? Our rule requires far less than trustworthy
metric pose. After Stage-3 alignment, a true re-edition presents the model
with two frames related by approximately the identity: every change that
reproduction introduces is photometric or planar-projective, there is no
parallax left to explain, and the pose head regresses near-coincident
cameras---near-zero $\delta_{\mathrm{rot}}$ and $\delta_{xy}$. A same-scene
negative---a second exposure taken steps away or moments later---cannot be
flattened this way: no single homography accounts for a genuine viewpoint
change of a non-planar scene, residual parallax survives alignment, and the
pose head absorbs it as relative camera motion. The signal is the
degeneracy itself: we use the pose head as a parallax detector, not a
camera localizer. Crucially, the near-zero pose response on true copies is
a measured behavior of the system, not an assumption \todo{synthetic
battery: $\sim$200 verified exact-copy pairs under controlled crops,
rotations, halftone simulation, and mild perspective warps; report the
fraction of pose scores remaining $\le 2.13$ per transform type --- run
W11} \todo{measured pose-score distributions per class on the labeled
shards, referenced from the Experiments section}.

\subsection{Design Alternatives}

The pose stage is one of several ways to extract viewpoint evidence from
signals the pipeline already computes, and we evaluate the alternatives
under an identical protocol on the same candidate sets. Decomposing the
Stage-3 homography yields a relative rotation and translation at no
additional cost \citep{malis2007homography}; the same correspondences
support an essential-matrix pose under assumed intrinsics
\citep{nister2004fivepoint}; MASt3R \citep{leroy2024mast3r} provides an
alternative learned pose head; and a small learned classifier over the
recorded per-pair signals (correspondence counts, $\rho$, the four pose
deltas, global similarity), in the spirit of Doppelgangers-style match
verification \citep{cai2023doppelgangers}, replaces the hand-set rule with
a fitted one. All four appear in Table~\ref{tab:verification}
\todo{Table~\ref{tab:verification} rows B7 (homography-decomposition pose),
B8 (essential-matrix pose), B9 (MASt3R pose, run W6), B10 (learned
classifier); B7/B8/B10 are CPU-only from cached manifests} alongside the
component and threshold ablations.

\section{Experiments}
\label{sec:evaluation}

Our evaluation answers six questions, each with its own denominator:
(1)~does the learned retrieval front-end outperform off-the-shelf
copy-detection and self-supervised descriptors on photo-to-halftone
retrieval (Table~\ref{tab:retrieval});
(2)~does the inferred-pose gate improve exposure-identity vs.\ same-scene
verification over cheaper geometric alternatives
(Table~\ref{tab:verification});
(3)~do the two hand-set thresholds sit on plateaus rather than spikes
(Figure~\ref{fig:sensitivity});
(4)~what fraction of known reproductions the full funnel recovers end to
end (Table~\ref{tab:survival});
(5)~does the frozen pipeline transfer to a second propaganda system; and
(6)~does the system measurably change historians' work.

\subsection{Protocol}
\label{sec:protocol}

\paragraph{Split roles.}
The working corpus is partitioned into six shards by deterministic
batching; three are labeled. \textbf{Shard~1 is the development set}: the
inlier-ratio cutoff (0.65), the pose-score threshold (2.13), and the
pose-component weighting itself were all derived here, and Shard-1 numbers
should be read as development-set results. \textbf{Shard~2 is held-out
validation}: it was evaluated once, with every parameter frozen, and never
used for tuning. \textbf{Shard~3 is an untouched test set}, labeled under
the written codebook of the Human Evaluation subsection and evaluated
exactly once, with everything frozen \todo{single-shot Shard-3 evaluation
(jobs W2/W3), scheduled after all thresholds and tables are locked}.
Labeled pairs: Shard~1, 325 positive / 316 negative (641); Shard~2, 252
positive / 386 negative (638); Shard~3 \todo{label counts after codebook
annotation}. Pairs labeled \emph{Unsure} are excluded from these counts and
reported as a first-class outcome, never silently dropped.

\paragraph{Denominator discipline.}
Every metric below names its denominator. \emph{Retrieval} Recall@$K$ and
mAP are computed over the 2{,}636 verified photo--print pairs, with the full
42{,}773-crop working corpus as distractors: they answer whether a true
reproduction, when one exists, surfaces in the top-$K$ list.
\emph{Verification} precision, recall, and F1 are \emph{conditional on the
candidate set}: ground-truth labels exist for the candidate pairs that
survive retrieval and the correspondence floor, and all metrics in
Table~\ref{tab:verification} are computed over these expert-adjudicated
pairs. We never quote a verification number as end-to-end performance.
The paper's \emph{only} end-to-end claim is the survival analysis of
Table~\ref{tab:survival}, whose denominator is again the 2{,}636 verified
pairs.

\paragraph{Statistics, seeds, compute.}
Every P/R/F1 and PR-AUC cell carries a bootstrap 95\% confidence interval
(10{,}000 resamples over pairs) \cite{efron1993bootstrap}
\todo{add CIs to all table cells once rows complete}. Variants are
compared to the full system with exact McNemar tests on paired per-pair
decisions \cite{mcnemar1947}, Holm-corrected within each table
\cite{holm1979}. The final representation model is retrained with three
seeds (mean$\pm$std) \todo{3-seed retrain, job W12}; ablation variants use
a single disclosed seed. All RANSAC estimation is seeded
(\texttt{cv2.setRNGSeed} before every call; fundamental-matrix RANSAC at a
1-pixel reprojection threshold) \cite{fischler1981ransac}, making the
geometric stages bit-for-bit reproducible; a five-repeat seed check on
Shard~1 confirms near-zero geometric-stage variance \todo{5-seed repeat
numbers}. The fusion head is trained transductively (self-supervised, on
the unlabeled evaluation corpus); a control retrain excluding all Shard-3
images is reported in \todo{transductive control, Shard-3 F1 within CI ---
run W9}. Experiments ran on \todo{GPU model(s), VRAM, total GPU-hours per
experiment family}. The decision layer, with both thresholds as named
constants, and scripts that regenerate every table from cached per-pair
manifests, are part of the anonymized code in the supplementary material.

\subsection{Retrieval: Candidate Proposal}

Table~\ref{tab:retrieval} compares the retrieval front-end against frozen
off-the-shelf descriptors under identical preprocessing, following the
distractor-corpus protocol of the Image Similarity Challenge
\cite{douze2021disc}. A GeM descriptor \cite{radenovic2019gem} represents
the classical landmark/instance-retrieval family (GeM/DELG) alongside the
copy-detection and foundation-model encoders. Two rows isolate our design
choices: DINOv3-alone with the SSL head (does the three-backbone fusion
beat its best member?) and fusion without SSL (does the SimCLR-style head
\cite{chen2020simclr} matter?).

\begin{table}[t]
\centering
\small
\begin{tabular}{@{}llccccc@{}}
\toprule
& Encoder & R@1 & R@5 & R@10 & R@15 & mAP \\
\midrule
A1 & SSCD \cite{pizzi2022sscd} & \multicolumn{5}{c}{\todo{W4}} \\
A2 & DINOv2 global \cite{oquab2024dinov2} & \multicolumn{5}{c}{\todo{W4}} \\
A3 & DINOv3 global \cite{simeoni2025dinov3} & \multicolumn{5}{c}{\todo{W4}} \\
A4 & CLIP / SigLIP~2 \cite{radford2021clip,tschannen2025siglip2} & \multicolumn{5}{c}{\todo{W4}} \\
A5 & GeM \cite{radenovic2019gem} & \multicolumn{5}{c}{\todo{W4}} \\
A6 & ViT ensemble \cite{du2024probing} & \multicolumn{5}{c}{\todo{W4}} \\
A7 & Ours: fusion + SSL & \multicolumn{5}{c}{\todo{W4}} \\
A8 & DINOv3-alone + SSL & \multicolumn{5}{c}{\todo{W9}} \\
A9 & Fusion, no SSL & \multicolumn{5}{c}{\todo{W9}} \\
\bottomrule
\end{tabular}
\caption{Retrieval over the 2{,}636 verified pairs with the 42{,}773-crop
corpus as distractors. \todo{All cells: W4 = featurize the corpus with
each frozen encoder and score Recall@$K$/mAP; W9 = retrain rows A8--A9
(DINOv3-alone + SSL head; fusion without SSL). A6 uses the ensemble
configuration of Du, Le, and Honig (2024), reimplemented from the
published description (confirm weights are publicly released before
claiming otherwise).}}
\label{tab:retrieval}
\end{table}

\subsection{Verification: Exposure Identity vs.\ Same Scene}

Table~\ref{tab:verification} is the paper's core table. All rows share
the identical, fixed candidate sets described in the Protocol subsection,
so differences are attributable to the decision rule alone; metrics are
conditional on the candidate set throughout. Row~B0 is the
accept-everything base rate, i.e.\ the class prior of the candidate set:
precision 50.7 on Shard~1 and 39.5 on Shard~2 at recall 100. Row~B2
(correspondence floor + inlier ratio, no pose) already reaches
82.9~/~97.2~/~89.5 (P~/~R~/~F1) on Shard~1 and 83.7~/~100~/~91.1 on
Shard~2. Adding the component-weighted pose gate (row~B3, the full
system) removes 17 false positives at the cost of 3 true positives on
Shard~1, and 22 false positives at the cost of 4 true positives on
Shard~2, yielding 86.7~/~96.3~/~91.3 and 90.2~/~98.4~/~94.1. Because the
pose gate only removes pairs from the inlier-gate survivors, the
discordant-pair counts for B2 vs.\ B3 are fully determined, and the exact
McNemar $p$-values are $2.6\times10^{-3}$ (Shard~1) and
$5.3\times10^{-4}$ (Shard~2); Holm correction is applied across the full
table once all rows are complete \todo{Holm-corrected $p$ column}.

Three groups of rows probe the mechanism. \emph{Component rows} B4/B5
split the pose score of Eq.~\ref{eq:pose-score} into its high-weight terms
(rotation angle and
$xy$-translation) and its low-weight terms (fov shift and
$z$-translation); if the weighting is right, B5 should perform near
chance, and the per-component score distributions
(Figure~\ref{fig:pose-components}) should show the same asymmetry. Row~B6
gates on the global feature-similarity signal that the pipeline records
but \emph{excludes} from its decision rule owing to heavy between-class
distribution overlap; B6 quantifies what a tuned gate on that signal would
have achieved. \emph{Alternative-pose rows} B7--B9 answer the question a
reader should ask: the alignment stage already estimates a homography, and
decomposing it yields relative rotation and translation for free
\cite{malis2007homography} (B7); an essential-matrix pose from the same
correspondences is nearly as cheap \cite{nister2004fivepoint} (B8); and
MASt3R \cite{leroy2024mast3r} is the strongest alternative feed-forward
pose head (B9). Each is tuned on Shard~1 only and applied frozen, exactly
as our gate was. Row~B10 replaces the hand-set thresholds with a learned
classifier over the stored per-pair signals, in the spirit of
Doppelgangers \cite{cai2023doppelgangers}. Rows B11--B14 perturb the
geometry stage itself (alignment off; filtered vs.\ raw alignment
correspondences; matcher swaps to LightGlue
\cite{lindenberger2023lightglue} and RoMa \cite{edstedt2024roma}).

\begin{figure}[t]
\centering
\fbox{\parbox{0.94\columnwidth}{\centering\vspace{1.5em}
\todo{per-component violin/ROC plots of the four pose deltas on exact-copy
vs.\ same-scene pairs, plus the global-similarity distribution-overlap
panel; CPU from \texttt{vggt\_judged\_manifest.jsonl}}
\vspace{1.5em}}}
\caption{Per-component pose-delta distributions for exposure-identical
vs.\ same-scene candidate pairs.}
\label{fig:pose-components}
\end{figure}

\begin{table*}[t]
\centering
\small
\begin{tabular}{@{}llccccc@{}}
\toprule
& & Shard 1 (dev) & Shard 2 (val) & Shard 3 (test) & PR-AUC & McNemar \\
& Variant & P~/~R~/~F1 & P~/~R~/~F1 & P~/~R~/~F1 & (S2) & vs.\ B3 (S2) \\
\midrule
B0 & Accept-all (base rate) & 50.7 / 100 / 67.3 & 39.5 / 100 / 56.6 & \todo{S3} & --- & $<10^{-6}$ \\
B1 & Correspondence floor only & \multicolumn{5}{l}{\todo{floor sweep 0--200 over \texttt{aspan\_all\_manifest} after breakpoint-0 rerun (W1); CPU}} \\
B2 & + inlier-ratio gate (no pose) & 82.9 / 97.2 / 89.5 & 83.7 / 100 / 91.1 & \todo{S3} & \todo{sweep} & $5.3\times10^{-4}$ \\
B3 & \textbf{+ weighted pose gate (full)} & \textbf{86.7 / 96.3 / 91.3} & \textbf{90.2 / 98.4 / 94.1} & \todo{S3} & \todo{sweep} & --- \\
\midrule
B4 & Pose: rotation + $xy$ terms only & \multicolumn{5}{l}{\todo{CPU re-scoring of cached manifests, \texttt{--pose-components rotation\_xy}}} \\
B5 & Pose: fov + $z$ terms only & \multicolumn{5}{l}{\todo{CPU re-scoring, \texttt{--pose-components fov\_z}; expected near-chance}} \\
B6 & Global-similarity gate & \multicolumn{5}{l}{\todo{threshold sweep on stored global-similarity signal; dev-tuned, frozen on S2/S3}} \\
\midrule
B7 & Homography-decomposition pose & \multicolumn{5}{l}{\todo{$R,t$ from \texttt{decomposeHomographyMat} on the cached alignment homographies; CPU}} \\
B8 & Essential-matrix pose & \multicolumn{5}{l}{\todo{five-point + RANSAC on cached ASpanFormer correspondences; CPU}} \\
B9 & MASt3R relative pose & \multicolumn{5}{l}{\todo{W6: MASt3R pose on all labeled pairs, same tune-on-S1 protocol}} \\
B10 & Learned classifier (LR / MLP) & \multicolumn{5}{l}{\todo{logistic reg.\ + 2-layer MLP over stored signals (correspondence count, inlier ratio, pose deltas, global sim); train S1, frozen}} \\
\midrule
B11 & Alignment off & \multicolumn{5}{l}{\todo{W8: VGGT on unaligned crops}} \\
B12 & Filtered alignment correspondences & \multicolumn{5}{l}{\todo{W8: raw vs.\ filtered correspondences for alignment}} \\
B13 & Matcher swap: LightGlue & \multicolumn{5}{l}{\todo{W7: replace ASpanFormer; re-tune both thresholds on S1, frozen on S2/S3}} \\
B14 & Matcher swap: RoMa & \multicolumn{5}{l}{\todo{W7: same protocol as B13}} \\
\bottomrule
\end{tabular}
\caption{Verification conditional on the candidate set: the expert-labeled
filter-survivor pairs of each shard (641 on Shard~1, 638 on Shard~2,
\todo{$n$} on Shard~3). Thresholds and weights for every variant are tuned
on Shard~1 only and applied frozen elsewhere. McNemar $p$-values are exact
tests on paired decisions; Holm correction across the table
\todo{apply once all rows are run}. Shard-3 column filled by the single
frozen evaluation \todo{W2/W3}.}
\label{tab:verification}
\end{table*}

\subsection{Threshold Sensitivity}

Both operating constants of Eq.~\ref{eq:decision} are hand-set, so we show
they lie on plateaus. Figure~\ref{fig:sensitivity} sweeps the inlier-ratio
cutoff over $[0.30, 0.90]$ and plots F1 per shard with 0.65 marked
(panel~a); gives full ROC and PR curves over the continuous pose score with
the 2.13 operating point marked (panel~b); and shows a two-dimensional F1
heatmap over (inlier ratio, pose score) on the development shard with the
chosen operating point, replotted on the untouched test shard to show that
the plateau, not just the point, transfers (panel~c). The frozen operating
point was evaluated on Shard~3 exactly once, before any sweep; the Shard-3
heatmap is a post-hoc robustness analysis computed after that single frozen
evaluation, with no feedback into any reported operating point.

\begin{figure}[t]
\centering
\fbox{\parbox{0.94\columnwidth}{\centering\vspace{1.5em}
\todo{sensitivity figure, three panels: (a) inlier-ratio sweep 0.30--0.90,
F1 per shard, 0.65 marked; (b) ROC/PR over the continuous pose score, 2.13
marked; (c) 2D F1 heatmap over (inlier ratio, pose score) on dev, replotted
on Shard 3 --- CPU sweeps over cached manifests; panel (c) requires Shard-3
manifests, W2/W3}
\vspace{1.5em}}}
\caption{Threshold sensitivity of the two frozen operating constants.}
\label{fig:sensitivity}
\end{figure}

\subsection{End-to-End Survival of Verified Pairs}

Verification metrics above say nothing about pairs the funnel never
surfaces. Table~\ref{tab:survival} therefore pushes all 2{,}636 verified
pairs through every stage at frozen thresholds and reports the fraction
still alive after each. This survival table is the paper's only
end-to-end claim, and its final row answers the historian's operative
question: \emph{of the reproductions known to exist, how many does the
system find?}

\begin{table}[t]
\centering
\small
\begin{tabular}{@{}lcc@{}}
\toprule
Stage & Pairs alive & \% of 2,636 \\
\midrule
All verified pairs & 2,636 & 100.0 \\
Retrieval top-10 & \todo{W5} & \todo{W5} \\
+ correspondence floor & \todo{W5} & \todo{W5} \\
+ inlier ratio $\geq 0.65$ & \todo{W5} & \todo{W5} \\
+ pose score $\leq 2.13$ (end-to-end) & \todo{W5} & \todo{W5} \\
\bottomrule
\end{tabular}
\caption{Stage-by-stage survival of the 2{,}636 verified photo--print
pairs at frozen thresholds \todo{W5: end-to-end recall audit; also report
which stage dominates the loss}.}
\label{tab:survival}
\end{table}

\subsection{Cross-Propaganda-System Transfer}

To test whether the pipeline and its thresholds encode anything specific
to one publication apparatus, we freeze everything --- weights,
thresholds, correspondence floor --- and run the identical system on the
Jinchaji Pictorial materials from the Papers of Sha Fei (2{,}146
photographs, 1{,}122 magazine pages, 4{,}390 crops, 230 annotated pairs):
Communist base-area propaganda, produced under different printing
technology and degradation than the Japanese-run \emph{North China}
corpus, and studied previously by \citet{du2024probing}, whose system
reached 77.8\% top-15 retrieval on this collection. We report Recall@10/15
against that reference \todo{W10: frozen-pipeline Jinchaji run} and
verification P/R/F1, conditional on the candidate set, on hand-labeled
filter survivors \todo{Jinchaji survivor labeling under the same
codebook}. For this collection the supplementary release contains
annotations and image identifiers only; the images themselves are not
redistributed.

\subsection{Human Evaluation}
\label{sec:human}

\paragraph{Annotation reliability.}
NCR-Match itself was built under the two-round verification protocol
described in Dataset Construction, whose second round excluded same-scene
negatives --- precisely the confusion the verification stage exists to
resolve. For the shard labels used in
Table~\ref{tab:verification}, a written one-page codebook fixes the
contested cases (adjacent contact-sheet frames; crops of a same-scene
negative), after which two annotators --- a domain-expert historian of the
period and a trained research assistant --- independently label all
candidate pairs as positive, negative, or \emph{Unsure}, without access
to each other's answers. We report Cohen's $\kappa$ \cite{cohen1960kappa}
\todo{$\kappa$ per shard from the dual relabeling}, per-shard
\emph{Unsure} counts \todo{counts}, and adjudicate disagreements into the
gold labels used throughout.

\paragraph{Historian time-saved study.}
We run a within-subjects study with 4--6 historians and graduate students
of the period \todo{final $N$ and recruitment description}. Each
participant works through two matched query sets of archival photographs
\todo{query-set size and matching procedure} under two conditions:
(A)~the full system --- pipeline-verified candidates with aligned overlay
visualizations --- and (B)~a plain top-10 retrieval list. Condition order
and query-set assignment are counterbalanced across participants. We
measure time-to-first-verified-republication (elapsed time until the
participant confirms a true republication of a query photograph) and
verified matches per 30-minute session. Because a participant-level
two-sided Wilcoxon signed-rank test \cite{wilcoxon1945} cannot reach
$p<0.05$ at this sample size (the minimum attainable $p$ is $0.125$ at
$N{=}4$ and $0.0625$ at $N{=}5$), the study is presented descriptively:
per-condition session medians, the direction of the paired difference
for every participant, and rank-biserial effect sizes, with sessions as
the named denominator and no hypothesis-test claim \todo{all results;
sessions scheduled}.

\subsection{Qualitative Analysis and Failure Modes}

Figure~\ref{fig:hard-cases} shows hard positives the system accepts
(heavy crops, coarse halftone screening, retouched reprints) alongside
hard negatives it rejects (same scene, exposures seconds apart), and
representative residual errors \todo{figure from adjudicated Shard-2
errors: 27 FP, 4 FN}. We additionally tag every confirmed positive with a
manipulation-trace type --- halftone-only, crop, rotation/mirror,
retouch, recaption-with-near-identical-pixels \cite{du2026evidence} ---
and report per-type pipeline pass rates \todo{taxonomy tagging of
confirmed positives; per-type survival}. One trace class is invisible to
this system by construction: recaption-only reuse, in which the printed
image is pixel-wise near-identical and only the surrounding caption
changes. The pipeline correctly verifies such pairs as the same exposure
--- that is its function --- but it cannot detect the caption
substitution itself. Detecting textual and layout-level traces around a
verified image pair is future work; the image-level verdicts reported
here are the evidentiary substrate such a system would require.

\begin{figure}[t]
\centering
\fbox{\parbox{0.94\columnwidth}{\centering\vspace{1.5em}
\todo{hard positives accepted (heavy crops, coarse halftone, retouch) vs.\
hard same-scene negatives rejected, plus representative residual errors
from the adjudicated Shard-2 error set (27 FP, 4 FN)}
\vspace{1.5em}}}
\caption{Hard cases and residual failure modes.}
\label{fig:hard-cases}
\end{figure}

\section{Findings for Historians}
\label{sec:findings}

The pipeline's product is not historical knowledge; it is a queue of
adjudicated photo--print pairs, each carrying its geometric evidence
(retrieval rank, inlier ratio, pose score, aligned overlay) and its
provenance metadata (publication, issue, date, page). It uses scale to
locate the singular object; the close reading of that object remains
the historian's work \cite{du2026evidence}. Findings of this kind have
precedent at proof-of-concept scale: \citet{du2024probing}
established the misattributed and mislocated Xifengkou/Futuyu
photograph and the four-printing recaptioning chain of Liu Hanxing's
enlistment portrait, and Du \citeyearpar{du2025metaphor} developed
their historiographic implications. We report here what the scaled
system adds: discoveries from the 42{,}773-crop working corpus that are,
to our knowledge, previously undocumented in the scholarship on either
archive, surfaced by the pipeline and adjudicated by a domain
historian.

\paragraph{Finding 1 (a new version chain).}
\todo{Insert a NEW version chain from the adjudicated accepted-pair
list: one North China Railway Archive photograph (give the archive
identifier) reproduced in two or more magazine placements with
materially different captions or layouts. Required elements: full
provenance line for every printing (publication = North China
(Hokushi), issue number, date, page); caption transcription and
translation; per-pair pipeline evidence (retrieval rank, inlier ratio,
pose score); one aligned-overlay figure. Must NOT be the
Xifengkou/Futuyu or Liu Hanxing cases (already published --- cited
above as prior findings). Source: domain-expert adjudication of the
Shard 1--3 accepted pairs, scheduled Jul 16--21.}

\paragraph{Finding 2 (a caption discrepancy or misattribution).}
\todo{Insert a NEW discrepancy between the archive's own metadata and
the printed caption (place, date, subject, or photographer credit),
with the same provenance-line and evidence requirements as Finding 1.
State the discrepancy only; assert no editorial intent --- intent is
the historian's argument, made elsewhere.}

\paragraph{Finding 3 (optional; an editorial crop or retouch).}
\todo{Insert only if adjudication yields it: a crop or retouch that
changes the legible content of the printed image (e.g., a person or
object removed by the crop), same evidence requirements. This slot is
cut first if space is short.}

These findings cost the historian hours, not months:
\todo{one sentence with the measured result of the within-subjects
historian study --- time-to-first-verified-republication and verified
matches per 30-minute session, pipeline-verified candidates with
overlays (condition A) vs.\ plain top-10 retrieval (condition B),
N=4--6 historians, 20--30 queries each; report medians and the paired
direction; denominator = sessions. Study of report section 6.6, sessions
Jul 16--21.}

\section{Ethics, Positionality, and Broader Impact}
\label{sec:ethics}

\paragraph{Two propaganda systems, one instrument.}
Both corpora are propaganda. The North China Railway Archive
photographs were produced by the press apparatus of a Japanese
corporate--military enterprise under occupation (1939--1945); the
Jinchaji material was produced by the Chinese Communist Party's
base-area pictorial unit, founded by the photographer Sha Fei. We
treat the two systems even-handedly, as parallel objects of forensic
description: the identical pipeline with identical frozen thresholds
is applied to both, and the paper makes no claim about the relative
truthfulness of either system. What it documents is a shared
mechanism --- the movement of photographs into recontextualized
print --- that both apparatuses operated.

\paragraph{Depicted persons and name restoration.}
The photographs show identifiable people living under occupation and
at war. Where circulation evidence shows an attribution was erased in
reprinting --- a photographer's credit dropped, a named subject
reduced to ``common people'' \cite{du2024probing} --- we follow a
name-restoration principle: re-attaching the documented name to the
image is archival justice and is stated as fact; any claim about
\emph{why} the name was removed is interpretation and is deferred to
historians.

\paragraph{The over-assertion risk.}
The system's core output --- ``these two images are the same
exposure'' --- is exactly the claim that, made falsely, launders a
misattribution into the scholarly record with a machine's authority.
Three design consequences follow. Decision thresholds are set
conservatively; every accepted pair is human-adjudicated before any
historical claim is made (verification precision, conditional on the
expert-labeled candidate set, is 90.2\% on the held-out validation
shard --- roughly one accepted pair in ten is wrong before
adjudication, which is why adjudication is not optional); and coverage
is reported as a per-stage survival table
(Table~\ref{tab:survival}, \todo{numbers from the end-to-end audit
over the 2{,}636 verified pairs, job W5}) rather than implied to be
complete.

\paragraph{Dual use.}
Nothing in the pose signal is specific to archives: a method that
separates copies from re-shots of the same scene transfers to
contemporary copy detection, misinformation triage, and copyright
enforcement, and, like all instance-matching technology, could serve
surveillance-adjacent uses. We release only archival material and
annotations, and we state the technique's generality rather than
leaving it implicit.

\paragraph{Dataset duties and licensing.}
The release ships with a datasheet \cite{gebru2021datasheets} that
records collection provenance, the two-round verification protocol,
and --- centrally --- the archive's own selection bias: the North
China Railway Archive is itself a curated product of the occupation
press apparatus, and what it omits is a property of the source that
the benchmark inherits; the benchmark measures matching, not
representativeness. North China Railway Archive images and metadata
are published under CC~BY~4.0 and are redistributed with attribution
to the archive \cite{ncrarchive2019}; for the Jinchaji side we release
annotations and image identifiers only, no images.

\paragraph{Design policy.}
The Introduction's design principle --- retrieval without
interpretation --- is implemented as policy, not aspiration: no output
field of the system asserts editorial intent, and no intent language
appears in the released annotations. Interpretation is reserved for
the historian, whose findings appear --- so labeled --- in the
preceding section.

\section{Limitations}
\label{sec:limitations}

\textbf{Image-only blindness.} The pipeline decides on pixels, and the
historically decisive manipulation is often textual: a recaption-only
reprint is correctly matched as the same exposure, but the discrepancy
that matters is invisible to the matcher, so caption changes are still
found only by reading every adjudicated pair.
\todo{Per-trace-type pipeline pass rates from the taxonomy tagging of
confirmed positives (halftone-only, crop, rotation/mirror, retouch,
recaption) --- quantifies exactly which manipulation classes the
image-only system cannot see.}
\textbf{Conditional metrics.} Verification P/R/F1 are conditional on
the candidate set; the survival table is the paper's only end-to-end
claim, and \todo{end-to-end survival numbers, audit W5} bound what a
historian ultimately misses.
\textbf{Transductive training.} The fusion head is trained
self-supervised on the full unlabeled working corpus, which includes
the evaluation images (no labels are used).
\todo{Control: retrain excluding Shard-3 images and show Shard-3 F1
unchanged within CI (report section 6.8, run W9); otherwise disclose as
an uncontrolled limitation.}
\textbf{Annotation reliability.} The shard labels originate from a
single annotation pass;
\todo{Cohen's kappa from the independent dual-annotation under the
written codebook, with disagreements adjudicated into gold labels and
Unsure as a first-class category.}
\textbf{Degenerate geometry by design.} An aligned exact copy is a
picture of a picture --- near-planar, near-zero baseline --- outside
the training distribution of a 3D reconstruction model, and the
near-zero inferred pose shift on true copies is a measured behavior,
\todo{synthetic battery: pose-score distributions under controlled
crops, rotations, halftone simulation, and mild perspective warps on
$\sim$200 exact-copy pairs (report section 6.5, job W11)}, not a
guaranteed property; it is validated for one pinned checkpoint only.
\textbf{Upstream crop dependency.} Every stage consumes YOLOv11 layout
crops \cite{khanam2024yolov11}; detection and crop-boundary errors
propagate silently.
\todo{Crop-stage QC: detection precision/recall on a 100-page audited
sample, plus one sentence on failure handling.}

\section{Conclusion}
\label{sec:conclusion}

We posed provenance verification for archival photographs as a
question the copy-detection literature names but does not decide with
geometry: same photograph, or same scene? Our contributions mirror
that framing. We formulated exposure-identity vs.\ same-scene
discrimination on cross-domain photo-to-halftone material and, to our
knowledge, are the first to decide it using the relative camera pose
inferred by a 3D foundation model, wrapped in a conservative,
auditable decision layer. We release NCR-Match --- 2{,}636 verified
photo--print pairs with exposure-identity and same-scene-negative
labels, with datasheet and code --- deposited, not promised. With
thresholds frozen on the development shard, the verifier reaches
P~90.2 / R~98.4 / F1~94.1 on held-out validation, conditional on the
candidate set (\todo{Shard-3 test row, one-shot frozen evaluation of
Jul 22}), and \todo{the Hokushi$\rightarrow$Jinchaji cross-system
transfer result, frozen thresholds, report section 6.7} carries the same
protocol across opposing propaganda systems. The findings of the
preceding section --- surfaced by the pipeline, adjudicated by a
historian --- are the intended unit of return. The pipeline retrieves;
the argument belongs to the historian.

\bibliography{refs}

% =============================================================================
% REPRODUCIBILITY CHECKLIST --- draft answers (trailing comment block).
% AAAI-27 mechanics: the kit's ReproducibilityChecklist.tex is dual-mode; it
% compiles standalone or is \input at the end of the paper before
% \end{document}. Answer by editing its 31 \question{...}{(yes/partial/no)}
% items. When ready:  % \input{authorkit/AuthorKit27/ReproducibilityChecklist}
% Draft answers below; items marked [contingent] flip to partial/no if the
% referenced \todo does not land in the July 31 supplementary.
%
% 1. GENERAL PAPER STRUCTURE
%   1.1 Conceptual outline/pseudocode of AI methods: YES (four-stage pipeline
%       description; decision rules stated with named constants).
%   1.2 Opinions/hypotheses delineated from facts: YES.
%   1.3 Pedagogical references for less-familiar readers: YES.
% 2. THEORETICAL CONTRIBUTIONS
%   2.1 Theoretical contributions: NO (2.2--2.8: NA).
% 3. DATASET USAGE
%   3.1 Relies on datasets: YES.
%   3.2 Motivation for selected datasets: YES.
%   3.3 Novel datasets included in a data appendix: YES [contingent on the
%       NCR-Match deposit shipping and kappa/gold-label adjudication
%       completing so the released labels are the adjudicated ones].
%   3.4 Novel datasets publicly available with a research license: YES ---
%       NCR side CC BY 4.0 with attribution; Jinchaji side annotations and
%       identifiers only [TODO: fix annotation license string].
%   3.5 Existing datasets cited: YES (NCR Archive; Jinchaji/Sha Fei materials
%       via Du, Le, and Honig 2024).
%   3.6 Existing datasets publicly available: PARTIAL --- NCR Archive is
%       public (CC BY 4.0); Papers of Sha Fei / Jinchaji Pictorial images are
%       not redistributable.
%   3.7 Non-public datasets described in detail with justification: YES.
% 4. COMPUTATIONAL EXPERIMENTS
%   4.1 Includes computational experiments: YES.
%   4.2 Hyperparameter ranges and selection criterion: PARTIAL --- 0.65/2.13
%       derived on Shard 1 (dev) and disclosed as such; becomes YES once the
%       sensitivity curves (inlier 0.30--0.90; pose ROC/PR; top-K 5--15) land.
%   4.3 Preprocessing code included: PARTIAL --- retrieval/filter/judge/
%       scoring code complete; [TODO: package YOLOv11 crop weights + training
%       annotations in the supplementary, else answer NO and disclose].
%   4.4 All experiment source code included: YES (anonymized code in
%       supplementary material).
%   4.5 Code publicly available with research license: YES [TODO: choose
%       repo license, e.g. MIT/Apache-2.0].
%   4.6 Code commented with references to paper steps: YES (per-step
%       docstrings; paper constants named in the decision layer).
%   4.7 Random seeds described sufficiently: PARTIAL --- geometric stage is
%       seeded (cv2 RANSAC, bit-for-bit); becomes YES after committing
%       torch/numpy/DataLoader seeding and the frozen-backbone eval-mode fix
%       (report section 6.0 item 3).
%   4.8 Computing infrastructure specified: PARTIAL --- becomes YES once the
%       compute paragraph (GPU types/hours, memory, OS, pinned library
%       versions from the environment file) is written.
%   4.9 Metrics formally described and motivated: YES (every metric names
%       its denominator: Recall@K over verified pairs; verification P/R/F1
%       conditional on the candidate set; end-to-end survival table).
%   4.10 Number of runs per result stated: YES --- single-seed ablations
%        disclosed; [contingent: 3-seed final model mean+-std (job W12),
%        else state 1 seed explicitly].
%   4.11 Analysis beyond point estimates: YES [contingent: bootstrap 95% CIs,
%        10,000 resamples over pairs, on every P/R/F1 cell].
%   4.12 Significance tests: YES [contingent: McNemar exact tests vs. full
%        pipeline, Holm-corrected].
%   4.13 All final hyperparameters listed: YES.
%
% Items whose answer depends on a \todo landing before July 31:
%   3.3, 3.4, 4.2, 4.3, 4.5, 4.7, 4.8, 4.10, 4.11, 4.12.
% =============================================================================

\end{document}
